The offset and the design make different contributions. Removing the offset from the exact same residual-strata samples increases random-split MAE, while leaving rankings exactly unchanged. Difficulty stratification and raw-response stratification are competitive score estimators; the raw-response variant even has slightly lower random and upper-tail point estimates than the proposed residual-centered construction. Thus the evidence supports source-informed random sampling more broadly, rather than a universal advantage for this particular embedding. Full Spearman results also show that raw-response strata are competitive for ranking (Table~\ref{tab:allrho}); the source-pair variance identity motivates a proxy rather than establishing a uniformly better embedding.

Variance-based allocation improves on proportional allocation in the random and upper-tail point estimates, but not in the lower tail. Source variance proxies can therefore remain imperfect even when target scores are estimated without design bias. The source-count diagnostic (Figure~\ref{fig:sources}) shows lower residual-strata point estimates with more available source models on the same seven larger cohorts, with wide benchmark-level uncertainty. Additional history helps in these settings but does not remove the observed budget or transfer limitations.
